\usepackage[utf8]{inputenc}
%slippery gypsy packages that need to come first, since others call them after
\usepackage[x11names,dvipsnames,svgnames,table]{xcolor}

% general incantations
\usepackage[export]{adjustbox}
\usepackage{afterpage}

\usepackage{graphicx}
\usepackage{placeins}
\usepackage{pdfpages}
\usepackage{algorithm2e}
\usepackage{array}
\usepackage{booktabs}
\usepackage[most]{tcolorbox}
\usepackage{calligra}
\usepackage{caption}
\usepackage{datetime}
\usepackage{dblfnote}
\usepackage{dirtytalk}
\usepackage{dsfont}
\usepackage{etex}
\usepackage{fancyhdr}
\usepackage{fix-cm}
\usepackage[T1]{fontenc}
\usepackage{textcomp,gensymb} %for \degree C symbol
\usepackage{graphicx}
\usepackage{lipsum}
\usepackage{listings}
\usepackage{transparent}
\usepackage[everyline=true,framemethod=tikz]{mdframed}
\usepackage{mparhack}
\usepackage{multicol}
\usepackage{multirow}
\usepackage{parskip}
\usepackage{lscape}
\usepackage{pdflscape}
\usepackage{pdfpages}
\usepackage{placeins}
\usepackage[document]{ragged2e}
\usepackage{rotating}
\usepackage{setspace}
\usepackage{subcaption}
\usepackage{threeparttable}
\usepackage[normalem]{ulem}
\usepackage{verbatim}
\usepackage{soul} %highlighting, strike through etc.

%Automated appendices
\usepackage[titletoc,title,header]{appendix} %advanced functionality

%language settings
\usepackage[utf8]{inputenc}
\usepackage[australian]{babel}
\usepackage{csquotes}

%page setup
%this where we adjust the binding offset, if relevant
\usepackage[a4paper,margin=2cm]{geometry}
\usepackage{lastpage} % for page 1 of n footers

%cross referencing
\usepackage[hidelinks]{hyperref}
\usepackage{cleveref}

% Quantum computing packages
\usepackage{tikz}
\usepackage{quantikz}

% Placeholder figure: draws a boxed, described stand-in for a figure not yet
% produced. Usage: \placeholderfigure[height]{caption}{label}{description}
\newcommand{\placeholderfigure}[4][6cm]{%
  \begin{figure}[htbp]
    \centering
    \fbox{%
      \begin{minipage}[c][#1][c]{0.86\linewidth}
        \centering\itshape\color{gray!60!black} #4
      \end{minipage}%
    }
    \caption{#2}
    \label{#3}
  \end{figure}%
}

% -----------------------------------------------------------------------------
%maths stuff
\usepackage{amsmath}
\numberwithin{equation}{section}

\usepackage{mathtools}
\usepackage{amssymb}
\usepackage{braket}

\usepackage{amsthm}
% Definitions usually use upright text
% Define a new theorem style for definitions
\newtheoremstyle{boldnote}
  {\topsep}       % Space above
  {\topsep}       % Space below
  {\normalfont}   % Body font (upright text for definitions)
  {}              % Indent amount
  {\bfseries}     % Theorem head font (this makes the header bold)
  {:}             % Punctuation after theorem head
  {.5em}          % Space after theorem head
  {\thmname{#1}\thmnumber{ #2}\thmnote{ (#3)}} % Theorem head spec

% 3. Apply the custom style to definitions
\theoremstyle{boldnote}
\newtheorem{definition}{Definition}

% 4. Apply the default plain style to theorems and corollaries
\theoremstyle{plain}
\newtheorem{theorem}{Theorem}
\newtheorem{proposition}{Proposition}
\newtheorem{corollary}{Corollary}

% -----------------------------------------------------------------------------
% Thesis notation (ThesisPlanning.md §8, decisions T-06 to T-09 and T-16).
% The Chebyshev operators are blackboard bold, as in Wu et al., under the macro
% names pihm's Sphinx configuration uses (\G, \Bt, \Dt), so an equation copies
% between the code documentation and the thesis unchanged. \G is the
% coefficient-space differentiation matrix, f -> f' (strictly upper
% triangular); Wu et al. print the same matrix as G_n^T.
\usepackage{bbm} % lowercase blackboard bold, for the product-to-sum fold
\newcommand{\G}{\mathbb{G}}               % differentiation matrix
\newcommand{\Bt}{\tilde{\mathbb{B}}}      % banded factor: \G = \Bt^{-1}\Dt
\newcommand{\Dt}{\tilde{\mathbb{D}}}      % banded factor
\newcommand{\M}{\mathbb{M}}               % multiplication: \M_{x^p} (lift), \M_a
\newcommand{\nfold}{\mathbbm{n}}          % product-to-sum fold: \nfold_1, \nfold_x
\newcommand{\Brow}{\mathbb{B}}            % rank-one zero (invariant) constraint row
\newcommand{\Dzero}{\mathbb{D}^{(0)}}     % rank-one collapse (regular) constraint row
\newcommand{\Ucg}{\mathbb{U}}             % Chebyshev--Gauss transform (nodal fold)
\newcommand{\Uodd}{U_{\mathrm{odd}}}      % average of the odd-offset truncated shifts
\newcommand{\Astd}{A_{\mathrm{std}}}      % standard-form residual
\newcommand{\Adesc}{A_{\mathrm{desc}}}    % descriptor-form residual
\newcommand{\Hstd}{H_{\mathrm{std}}}      % standard-form Hamiltonian
\newcommand{\Hdesc}{H_{\mathrm{desc}}}    % descriptor-form Hamiltonian
\newcommand{\fhat}[1]{\hat{f}^{(#1)}}     % coefficients of the #1-th derivative (a descriptor block)

% -----------------------------------------------------------------------------

\setcounter{secnumdepth}{3}

%lists
\usepackage{enumitem}

%working collaboratively
\usepackage[backgroundcolor=yellow]{todonotes}

% Provisional numbers (decision T-15): pihm's numbers are quoted with \res
% (0_results/results.tex, input by main.tex). \prov{...} wraps a number pihm
% does not produce yet, and renders coloured until a \res replaces it. The
% pre-submission check (ThesisPlanning.md §14) fails on any remaining \prov.
\newcommand{\prov}[1]{\textcolor{BrickRed}{#1}}

%reference list 
\usepackage[backend=biber,style=ieee]{biblatex} %Sbliography with IEEE style.
\addbibresource{3_footer/ref.bib} %Imports bibliography file
\usepackage{url} %for URL fonts
\urlstyle{same}

%glossary for acronyms
\usepackage[acronym,nonumberlist,toc,section=subsection,numberedsection=nolabel]{glossaries} 
\makeglossaries
\newacronym{qsp}{QSP}{quantum signal processing}
\newacronym{qsvt}{QSVT}{quantum singular value transformation}
\newacronym{gqsp}{GQSP}{generalised quantum signal processing}
\newacronym{qite}{QITE}{quantum imaginary-time evolution}
\newacronym{oaa}{OAA}{oblivious amplitude amplification}
\newacronym{nisq}{NISQ}{noisy intermediate-scale quantum}
\newacronym{de}{DE}{differential equation}
\newacronym{ode}{ODE}{ordinary differential equation}
\newacronym{pde}{PDE}{partial differential equation}
\newacronym{nde}{NDE}{nonlinear differential equation}
\newacronym{dae}{DAE}{differential-algebraic equation}
\newacronym{pihm}{PIHM}{physics-informed effective Hamiltonian method}
\newacronym{hhl}{HHL}{Harrow--Hassidim--Lloyd}
\newacronym{lcu}{LCU}{linear combination of unitaries}
\newacronym{nlft}{NLFT}{nonlinear Fourier transform}
\newacronym{ir}{IR}{intermediate representation}
\newacronym{mcx}{MCX}{multi-controlled $X$}
\newacronym{qft}{QFT}{quantum Fourier transform}
\newacronym{dct}{DCT}{discrete cosine transform}
\newacronym{svd}{SVD}{singular value decomposition}
\newacronym{ivp}{IVP}{initial-value problem}
\newacronym{bvp}{BVP}{boundary-value problem}
\newacronym{dns}{DNS}{direct numerical simulation}

%line spacing
\linespread{1.25}